\documentclass[10pt,a4paper]{amsart}
\usepackage{amsmath,amssymb,amsthm,graphicx}
\RequirePackage{amssymb,amsmath,amsthm,bm}
\RequirePackage{thmtools}
\RequirePackage[capitalize]{cleveref}
\providecommand{\tightlist}{%
  \setlength{\itemsep}{0pt}\setlength{\parskip}{0pt}}

%%%%% varphi instead of phi
\renewcommand{\phi}{\varphi}

\mathchardef\mhyph="2D 
\providecommand{\A}{\mathcal{A}}
\providecommand{\C}{\mathcal{C}}
\providecommand{\L}{\mathcal{L}}
\providecommand{\B}{\mathcal{B}}
\providecommand{\M}{\mathcal{M}}
\renewcommand{\S}{\mathcal{S}}
\providecommand{\ol}[1]{\overline{#1}}
\providecommand{\mbf}[1]{\mathbf{#1}}
\providecommand{\mf}[1]{\mathfrak{#1}}
\providecommand{\mrm}[1]{\mathrm{#1}}
\providecommand{\mc}[1]{\mathcal{#1}}
\providecommand{\ra}{\rightarrow}
\providecommand{\Ra}{\Rightarrow}
\providecommand{\La}{\Leftarrow}
\providecommand{\lr}{\leftrightarrow}
\providecommand{\hf}{\mathbb{HF}}
\providecommand{\spec}{DgSp}
\providecommand{\bSigma}{\bm\Sigma}
\providecommand{\bDelta}{\bm\Delta}
\providecommand{\bPi}{\bm\Pi}
\providecommand{\Pinf}[1]{\Pi^{\mathrm{in}}_{#1}}
\providecommand{\Sinf}[1]{\Sigma^{\mathrm{in}}_{#1}}
\providecommand{\dSinf}[1]{d\mhyph\Sigma^{\mathrm{in}}_{#1}}
\providecommand{\Dinf}[1]{\Delta^{\mathrm{in}}_{#1}}
\providecommand{\bp}{\bar p}
\providecommand{\bq}{\bar q}
\providecommand{\ba}{\bar a}
\providecommand{\bb}{\bar b}
\providecommand{\bc}{\bar c}
\providecommand{\bd}{\bar d}
\providecommand{\bx}{\bar x}
\providecommand{\by}{\bar y}
\providecommand{\bz}{\bar z}
\providecommand{\LO}{\mathrm{LO}}
\declaretheorem[name={Theorem},numberwithin=section]{theorem}
\declaretheorem[name={Proposition}, sibling=theorem]{proposition}
\declaretheorem[name={Lemma}, sibling=theorem]{lemma}
\declaretheorem[name={Fact}, sibling=theorem]{fact}
\declaretheorem[name={Claim}, numberwithin=theorem]{claim}
\declaretheorem[name={Definition},numberwithin=section,sibling=theorem, style=definition]{definition}
\declaretheorem[name={Remark}, numbered=no,style=definition]{remark}
\declaretheorem[name={Corollary}, sibling=theorem]{corollary}
\declaretheorem[name={Question},sibling=theorem]{question}
\newcommand{\T}{{\mathcal{T}}}
\newcommand{\SR}{SR}
\newcommand{\pSR}{pSR}
\newcommand{\Q}{\mathbb{Q}}
\newcommand{\Z}{\mathbb{Z}}
\newcommand{\K}{\mathcal{K}}
\newcommand{\N}{\mathcal{N}}
\renewcommand{\L}{\mathcal{L}}
\providecommand{\Ginf}[1]{\Gamma^{\mathrm{in}}_{#1}}
\title{Appending an incomparable tuple does not change the back-and-forth relation of tuples}
\author{David Gonzalez, Dino Rossegger}
\date{arxiv:2305.07126v1; CC BY 4.0}
\begin{document}
\maketitle
\noindent\textit{Editorial note.} The counted statement is the author environment at source0.tex lines 346-349 and its proof environment opens at line 350, the line immediately after; the pairing is by adjacency and the retained proof range 350-367 is the whole environment.\par\medskip
\noindent\textit{Editorial note.} The counted statement uses the notions of somewhat comparable and entirely incomparable points and tuples, which the author introduces immediately before it at source0.tex lines 341-345; those definitions are retained as a context span (the same source, so the reader sees the author wording rather than a paraphrase), and they are not counted.\par\medskip
\noindent\textit{Editorial note.} The retained preamble is the author style file dino.sty (package lines, the theorem declaration support, the notation and tuple macros and the theorem environments) followed by the ordinal and rank macros of the manuscript itself, reproduced verbatim; this is the same recipe as the previously published card of this article, so no package outside it is requested and extra\_packages stays empty.\par\medskip
\noindent\textit{Editorial note.} Metadata provenance: source\_id and source\_version are the values recorded in this source licence evidence file (sources/scott-linear-orders-2305.07126v1/LICENSE\_EVIDENCE.json: arxiv:2305.07126v1 / 2305.07126v1), which the card points at directly; the title and author list are transcribed from the authors\_index and title keys of the source record of the pinned copy. The card carries no counter setting and imposes no original numbering, since this article has no label-number dump and its numbering cannot be reconstructed from the retained ranges alone.\par\medskip
\noindent\textit{Editorial note.} No external reference is needed: the retained statement, the definitions in context and the proof use no label or citation from outside the retained ranges.\par\medskip
\noindent\textit{Editorial scope.} One original result of the pinned author manuscript, counted once: appending a tuple incomparable to both sides preserves the back-and-forth relation of two tuples (author label lem:bftos\_noparameters).\par\medskip
\noindent\textit{Allowed standard prerequisites (editorial).}\begin{itemize}
\item Tree of tuples. The labelled tree T\_A of a structure is the tree of its finite tuples ordered by extension.. It is the tree on which the tuples live.
\item Somewhat and entirely incomparable points. A point is somewhat comparable to a tuple c when some non-root prefix of the point is a prefix of some entry of c, and entirely incomparable when it is not; a tuple is entirely incomparable when all of its points are.. It is the hypothesis on the two tuples and is retained inside the card as context.
\item Back-and-forth relations of the tree. The relations <=\_alpha are the back-and-forth relations of the labelled tree, defined by the usual game of rank alpha.. They are the relation the lemma compares.
\end{itemize}\medskip
say that $y\in T$ is \textit{somewhat comparable} to $\bar{c}$ if $\exists z,i ~
z\neq r\land c_i\succeq z \land y\succeq z$.
A tuple is somewhat comparable to $\bar{c}$ if all of its points are.
Similarly, we say that a point is \textit{entirely incomparable} to $\bar{c}$ if it is not somewhat comparable to $\bar{c}$ and that a tuple is entirely incomparable to $\bar{c}$ if all of its points are entirely incomparable to $\bar{c}$.


\begin{lemma}\label{lem:bftos_noparameters}
For $\T_\A$ a tree of tuples, every $\alpha\in\omega_1$ and tuples $\ba$ and $\bb$ that are entirely incomparable to $\bp$,
\[\ba\bp\leq_\alpha\bb\bp \iff \ba\leq_\alpha\bb.\]
\end{lemma}

\begin{proof}
We go by transfinite induction on $\alpha$.
The $\alpha=0$ case is immediate, as is the limit case.
Assume the result for $\alpha=\beta$.
If $\ba\bp\leq_{\beta+1}\bb\bp$, then $\ba\leq_{\beta+1}\bb$ follows immediately.
Now assume that $\ba\leq_{\beta+1}\bb$.
We wish to show that $\ba\bp\leq_{\beta+1}\bb\bp$.
Consider some extension $\bc$.
We can split $\bc$ into the points somewhat comparable to $\bp$, named $\bc_p$, and the points entirely incomparable to $\bp$, named $\bc_i$ (one of these may be the empty tuple).
Because $\ba\leq_{\beta+1}\bb$, there is a tuple $\bd_i$ such that $\ba\bd_i\geq_{\beta}\bb\bc_i$.
Without loss of generality (using an automorphism that fixes $\ba$ and moves $\bp$ to something entirely incomparable if needed), $\bd_i$ is entirely incomparable to $\bp$. 
Indeed, as somewhat comparability is an equivalence relation, we can note that $\ba\bd_i$ and $\bb\bc_i$ are entirely incomparable with $\bp\bc_p$.
In particular, the induction hypothesis gives that
\[\ba\bd_i\bp\bc_p\geq_{\beta}\bb\bc_i\bp\bc_p.\]
Rearranging the tuples gives that
\[\ba\bp\bd_i\bc_p\geq_{\beta}\bb\bp\bc\bc_p\]
which shows that $\ba\bp\leq_{\beta+1}\bb\bp$, as desired.
\end{proof}

\end{document}
